\documentclass[12pt,a4paper]{article}
\usepackage[dvipsnames]{xcolor}
\usepackage[utf8]{inputenc}
\usepackage{amsmath}
\usepackage{mathtools}
\usepackage{marvosym}
\usepackage{wrapfig}
\usepackage{hyperref}
\usepackage{float}
\usepackage{multicol}
\hypersetup{colorlinks,citecolor=black,filecolor=black,linkcolor=black,urlcolor=black}
\usepackage{pdfpages}
\usepackage{amsfonts}
\usepackage{amssymb}
\usepackage{fancyhdr}
\usepackage{chemfig}
\usepackage{graphicx}
\usepackage[T1]{fontenc}
\usepackage[magyar]{babel}
\usepackage{bm}
\usepackage{tikz, tcolorbox}
\usetikzlibrary{positioning}
\usepackage{verbatim}
\usepackage{amsmath,esint}
\usepackage{setspace}
\usepackage{qtree}
\usepackage{multido}
\usepackage{mathptmx} % Times New Roman betűtípus
\usepackage{listings}
\usepackage{caption}
\captionsetup[lstlisting]{labelformat=empty}

\lstset{
    language=Matlab,
    inputencoding=utf8,
    extendedchars=true,
    basicstyle=\ttfamily\small,
    frame=single,
    breaklines=true,
    numbers=left,
    numberstyle=\tiny,
    literate=
        {á}{{\'a}}1
        {é}{{\'e}}1
        {í}{{\'i}}1
        {ó}{{\'o}}1
        {ö}{{\"o}}1
        {ő}{{\H{o}}}1
        {ú}{{\'u}}1
        {ü}{{\"u}}1
        {ű}{{\H{u}}}1
        {Á}{{\'A}}1
        {É}{{\'E}}1
        {Í}{{\'I}}1
        {Ó}{{\'O}}1
        {Ö}{{\"O}}1
        {Ő}{{\H{O}}}1
        {Ú}{{\'U}}1
        {Ü}{{\"U}}1
        {Ű}{{\H{U}}}1
}
\renewcommand{\figureautorefname}{ábra}
\newcommand{\abrara}[1]{\hyperref[#1]{\ref*{#1}.~ábra}}
\newcommand{\Pointilles}[1]{%
\par\nobreak
\noindent\rule{0pt}{1.5\baselineskip}% Provides a larger gap between the
\multido{}{#1}{\noindent\makebox[\linewidth]{\dotfill}\endgraf}
\bigskip% Gap between dots and next paragraph
}
\newcommand{\feladat}[1]{%
\phantomsection
\subsection*{#1. feladat}%
\addcontentsline{toc}{subsection}{#1. feladat}%
}
\usepackage{pgfplots}
\pgfplotsset{height = 10cm, width=15cm,compat=1.9}
\usepackage[left=2cm,right=2cm,top=2cm,bottom=2cm]{geometry}
\setlength{\parindent}{0pt}
\setlength{\parskip}{0em}
\pagestyle{fancy}
\fancyhf{}
\cfoot{\thepage. oldal}

% Másfeles sorköz beállítása
\onehalfspacing

\begin{document}

\begin{titlepage}
\centering

\includegraphics[width=.8\textwidth]{../../Kepek/bme_logo_nagy.jpg}

\vspace{1em}

{
\large
BUDAPESTI MŰSZAKI ÉS GAZDASÁGTUDOMÁNYI EGYETEM\\ VILLAMOSMÉRNÖKI ÉS INFORMATIKAI KAR\\ MESTERSÉGES INTELLIGENCIA ÉS RENDSZERTERVEZÉS TANSZÉK
}

\vspace{6em}

{
\Huge
\textbf{Digitális szűrők tervezése és implementációja}
}

\vspace{3em}

{
\huge
Intelligens beágyazott rendszerek laboratórium\\

\vspace{.2em}

\large
(BMEVIMIMA21)
}

\vspace{9em}

{
\Large
\textit{Készítette:}\\
\vspace{.5em}
\textbf{Pál Patrik, Gavallér Dávid}\\
\vspace{.5em}
\textbf{AIQTQW, ES02RO}
}

\vspace{12em}

{
\large
BUDAPEST, \today
}
\end{titlepage}

\tableofcontents

\newpage

\section{Bevezetés}

\subsection*{A mérés célja}

A mérés során FIR és IIR digitális szűrőket vizsgáltunk. Az első feladatban egy megadott frekvenciahatárokhoz tartozó FIR aluláteresztő szűrőt terveztünk MATLAB-ban. A második feladatban a megadott együtthatókészletből IIR szűrőt készítettünk, és megnéztük annak átviteli függvényét.

\subsection*{A méréshez használt eszközök}
\begin{itemize}
    \item Számítógép
    \item MATLAB
    \item CrossCore Embedded Studio
    \item ADSP-BF537 EZ-KIT Lite fejlesztőpanel
    \item Digitális függvénygenerátor
    \item Oszcilloszkóp FFT funkcióval
    \item Hangforrás és hangsugárzó
\end{itemize}

\newpage

\section{Mérési feladatok}

\feladat{1}

\subsubsection*{Feladat leírása}

Tervezzen FIR aluláteresztő szűrőt 5~kHz-es határfrekvenciával és 6~kHz-es stopfrekvenciával. A passzsávi hullámosság legyen 2~dB, a stop-sáv csillapítása pedig legalább 30~dB. A szűrő együtthatóit töltse be a Blackfin projektbe, majd söpört szinusz jellel és az oszcilloszkóp FFT-jével ellenőrizze a szűrőt.

\subsubsection*{Feladat megoldása}

Első lépésben a Blackfin panelt és az oszcilloszkópot kötöttük össze. A mérési összeállítás a(z) \abrara{fig:meresi-elrendezes}-n látható.

\begin{figure}[H]
    \centering
    \includegraphics[width=0.52\textwidth]{../../Kepek/20261019_165007.jpg}
    \caption{A Blackfin panel és az oszcilloszkóp a mérési összeállításban.}
    \label{fig:meresi-elrendezes}
\end{figure}

A MATLAB-kódban először megadtuk a mintavételi frekvenciát és a szűrő határfrekvenciáit. A \texttt{firpmord} függvény megadta a tervezéshez szükséges szűrőfokot és súlyokat, az együtthatókat pedig a \texttt{firpm} függvénnyel számítottuk ki. Az átvitelt a \texttt{freqz} függvénnyel ábrázoltuk. Az ábrán a piros vonalak jelölik az előírt határértékeket.

Az együtthatókat a \texttt{FIR\_lowpass\_coef.dat} fájlba írtuk, ahonnan a CrossCore projekt FIR ága beolvassa azokat. A számított átviteli függvény a(z) \abrara{fig:fir-atvitel} látható.

\begin{figure}[H]
    \centering
    \includegraphics[width=\textwidth]{../../Kepek/FIR_atvitel.jpg}
    \caption{A FIR aluláteresztő szűrő átviteli függvénye.}
    \label{fig:fir-atvitel}
\end{figure}

A Blackfin kimeneti jelét az oszcilloszkóp FFT funkciójával is megvizsgáltuk, ezt a(z) \abrara{fig:fir-meres} szemlélteti.

\begin{figure}[H]
    \centering
    \includegraphics[width=0.9\textwidth]{../../Kepek/20261019_165015.jpg}
    \caption{A FIR szűrő kimenetének oszcilloszkópos FFT-képe.}
    \label{fig:fir-meres}
\end{figure}

\newpage

\feladat{2}

\subsubsection*{Feladat leírása}

Készítsen IIR szűrő rutint a Blackfin FIR projekt alapján. A rutin tetszőleges rendű együtthatók kezelésére legyen alkalmas. Használja a mérési útmutatóban megadott együtthatókészletek egyikét.

\subsubsection*{Feladat megoldása}

A 4. együtthatókészletet használtuk. A szűrő különbségegyenlete:

\begin{equation}
    y[n] = \sum_{k=0}^{4} b_k x[n-k] - \sum_{k=1}^{4} a_k y[n-k].
\end{equation}

A kiválasztott együtthatók közelítő értékei:

\begin{equation}
    \begin{aligned}
        \boldsymbol{b} &= [0{,}4998,\;-1{,}9993,\;2{,}9989,\;-1{,}9993,\;0{,}4998],\\
        \boldsymbol{a} &= [1,\;-2{,}6386,\;2{,}7693,\;-1{,}3393,\;0{,}2498].
    \end{aligned}
\end{equation}

A Blackfin programban az együtthatókat néggyel osztottuk, mert a nagyobb értékek nem férnek el a 16 bites \texttt{fract} formátumban. A szűrés végén az eredményt néggyel szoroztuk. A bemeneti és kimeneti mintákat külön tömbökben tároltuk, és a ciklusokkal összegeztük a megfelelő mintákat. A mérést az előző feladatnál bemutatott összeállításon végeztük.

A Blackfin kimenetét az oszcilloszkóp FFT funkciójával ismét megvizsgáltuk. A(z) \abrara{fig:iir-meres}-n kis frekvencián alacsonyabb, nagyobb frekvencián közel állandó kimeneti szint látható, ez felüláteresztő szűrőre utal.

\begin{figure}[H]
    \centering
    \includegraphics[width=0.75\textwidth]{../../Kepek/20261019_173704.jpg}
    \caption{A 4. IIR együtthatókészlet oszcilloszkópos FFT-képe.}
    \label{fig:iir-meres}
\end{figure}


\newpage
\feladat{3}

\subsubsection*{Feladat leírása}

Tervezzen 4. rendű Butterworth IIR aluláteresztő szűrőt a \texttt{butter} MATLAB-függvénnyel. A határfrekvencia legyen 12~kHz.

\subsubsection*{Feladat megoldása}

A tervezéshez 48~kHz-es mintavételi frekvenciát használtunk. A MATLAB-ban a határfrekvenciát a Nyquist-frekvenciával osztottuk, ezért a \texttt{butter} függvénynek $12000/(48000/2)=0{,}5$ normalizált frekvenciát adtunk meg. A kapott együtthatók:

\begin{equation}
    \begin{aligned}
        \boldsymbol{b} &= [0{,}093981,\;0{,}375923,\;0{,}563885,\;0{,}375923,\;0{,}093981],\\
        \boldsymbol{a} &= [1,\;0,\;0{,}486029,\;0,\;0{,}017665].
    \end{aligned}
\end{equation}

A visszacsatoló együtthatók és a számláló együtthatói is elférnek a \texttt{fract} formátumban, ezért ennél a szűrőnél nem kellett osztást alkalmazni. Az átviteli függvény a(z) \abrara{fig:butterworth-atvitel} ábrán látható. A határfrekvencián az amplitúdó körülbelül $-3$~dB.

\begin{figure}[H]
    \centering
    \begin{tikzpicture}
    \begin{axis}[
        width=0.9\textwidth,
        height=8cm,
        xlabel={Frekvencia [Hz]},
        ylabel={Amplitúdó [dB]},
        xmin=0, xmax=24000,
        ymin=-80, ymax=5,
        grid=both
    ]
    \addplot table [x=f, y=db] {../../Kepek/Butterworth_atvitel.dat};
    \end{axis}
    \end{tikzpicture}
    \caption{A Butterworth szűrő átviteli függvénye.}
    \label{fig:butterworth-atvitel}
\end{figure}

A szűrő a mérésvezetővel megbeszélve, nem lett implementálva a Blackfin panelen.


\newpage
\feladat{4}

\subsubsection*{Feladat leírása}

Tervezzen 4. rendű Csebisev I IIR aluláteresztő szűrőt a \texttt{cheby1} MATLAB-függvénnyel. A határfrekvencia legyen 4~kHz, a passzsávi hullámosság pedig 2~dB. Hasonlítsa össze a szűrőt a Butterworth szűrővel!

\subsubsection*{Feladat megoldása}

A mintavételi frekvencia most is 48~kHz volt. A \texttt{cheby1} függvény bemenetén a 4~kHz-es határfrekvenciát $4000/(48000/2)=0{,}1667$ értékkel adtuk meg, a passzsávi hullámosság pedig 2~dB lett. Az együtthatók:

\begin{equation}
    \begin{aligned}
        \boldsymbol{b} &= [0{,}000652,\;0{,}002606,\;0{,}003909,\;0{,}002606,\;0{,}000652],\\
        \boldsymbol{a} &= [1,\;-3{,}371443,\;4{,}505325,\;-2{,}808567,\;0{,}687808].
    \end{aligned}
\end{equation}

A passzsávban a válasz 0 és $-2$~dB között változik, látszik a megadott hullámosság, a Butterworth szűrő passzsávja viszont sima. A Csebisev I szűrő határfrekvenciája után a görbe meredekebben csökken. A két szűrő határfrekvenciája eltér, ezért az összehasonlító ábrán a frekvenciát a saját határfrekvenciához viszonyítottuk. A Csebisev I átviteli függvénye a(z) \abrara{fig:cheby1-atvitel}, az összehasonlítás a(z) \abrara{fig:iir-osszehasonlitas}-n látható.

\begin{figure}[H]
    \centering
    \begin{tikzpicture}
    \begin{axis}[
        width=0.9\textwidth,
        height=8cm,
        xlabel={Frekvencia [Hz]},
        ylabel={Amplitúdó [dB]},
        xmin=0, xmax=24000,
        ymin=-80, ymax=5,
        grid=both
    ]
    \addplot table [x=f, y=db] {../../Kepek/Cheby1_atvitel.dat};
    \end{axis}
    \end{tikzpicture}
    \caption{A Csebisev I szűrő átviteli függvénye.}
    \label{fig:cheby1-atvitel}
\end{figure}

\begin{figure}[H]
    \centering
    \begin{tikzpicture}
    \begin{axis}[
        width=0.9\textwidth,
        height=8cm,
        xlabel={Normalizált frekvencia ($f/f_c$)},
        ylabel={Amplitúdó [dB]},
        xmin=0, xmax=2,
        ymin=-80, ymax=5,
        grid=both,
        legend pos=south west
    ]
    \addplot table [x=fn, y=butter] {../../Kepek/Butterworth_Cheby1_osszehasonlitas.dat};
    \addplot table [x=fn, y=cheby] {../../Kepek/Butterworth_Cheby1_osszehasonlitas.dat};
    \addplot[dashed, black] coordinates {(1,-80) (1,5)};
    \legend{Butterworth, Csebisev I, $f/f_c=1$}
    \end{axis}
    \end{tikzpicture}
    \caption{A Butterworth és a Csebisev I szűrő összehasonlítása, a saját határfrekvenciájukra normálva.}
    \label{fig:iir-osszehasonlitas}
\end{figure}

A szűrő a mérésvezetővel megbeszélve, nem lett implementálva a Blackfin panelen.

\newpage
\section*{Függelék}
\subsection*{MATLAB kódok}

\lstinputlisting[
    caption={A FIR szűrő MATLAB-kódja},
    label={lst:fir-matlab}
]{../../MATLAB/feladat1.m}

\lstinputlisting[
    caption={A Butterworth szűrő MATLAB-kódja},
    label={lst:butterworth-matlab}
]{../../MATLAB/feladat3.m}

\lstinputlisting[
    caption={A Csebisev I szűrő MATLAB-kódja},
    label={lst:cheby1-matlab}
]{../../MATLAB/feladat4.m}

\end{document}
